% !TEX root = ../../LAFROX_AnalisisFinanciero_2_20261113.tex
\section{Modelo de costos operacionales}
\label{sec:e2-costos-operacionales}

Los 36 meses de Operación, del 21 al 56, cuestan 109.041,65 UF netos, o \$4.361.665.946 y USD~4.846.295. En promedio son 3.028,93 UF al mes, frente a una cuota de 3.550,01 UF que paga el CLIENTE. Dos tercios del costo son personas, y la mayor parte del resto son servicios de nube y comunicaciones.

Las cifras salen del modelo financiero (LafroX, 2026f), en sus hojas \texttt{COSTOS}, \texttt{COSTOS HH - SUELDOS}, \texttt{ADQUISICIONES OPERACION} y \texttt{GASTOS ADMINISTRATIVOS}. Las cantidades y los horarios de servicio vienen del T-11, de las cuentas 8.1, 8.2, 8.4 y 8.5 del T-14 y de los perfiles del T-15 (LafroX, 2026a, 2026c).

\subsection{Composición del presupuesto y reglas de imputación}
\label{sec:e2-op-presupuesto}

La Tabla~\ref{tab:e2-op-presupuesto} cuadra con el total de soporte de la hoja \texttt{RESUMEN} y considera solo lo que se paga entre los meses 21 y 56. Lo comprado en la implementación no se vuelve a cargar.

\begin{tablalafrox}{Costo de Operación de los meses 21 a 56, neto de IVA. Fuente: modelo financiero, hojas COSTOS y RESUMEN.}{tab:e2-op-presupuesto}{Y R{2.2cm} R{3.0cm} R{2.4cm}}
  {\cab{Categoría} & \cab{UF} & \cab{CLP} & \cab{USD}}
  Recursos humanos & 74.273,19 & 2.970.927.600 & 3.301.031 \\
  Capacitación, mentoría y certificación & 1.933,53 & 77.341.020 & 85.934 \\
  Infraestructura y equipamiento & 671,66 & 26.866.393 & 29.852 \\
  Servicios de proveedores & 25.623,76 & 1.024.950.443 & 1.138.834 \\
  Gastos administrativos & 6.225,30 & 249.012.090 & 276.680 \\
  Reservas de riesgo & 314,21 & 12.568.400 & 13.965 \\
  \midrule
  \textbf{Total de Operación} & \textbf{109.041,65} & \textbf{4.361.665.946} & \textbf{4.846.295} \\
\end{tablalafrox}

El costo cambia de un mes a otro. Las remuneraciones siguen las HH programadas, la reposición de equipos cae en el mes 21 y las garantías, los seguros y la custodia de fuentes se renuevan una vez al año (ver el flujo mensual de la sección~\ref{sec:e2-curva-s}). Los intereses y la devolución del crédito van en el financiamiento.

Ningún costo se cuenta dos veces. Las HH de capacitación van separadas de las del servicio, salvo la formación de los operadores del NOC para la franja de mesa, que es parte de sus horas, la nube de las innovaciones está dentro de los servicios de proveedores, y garantías y seguros están dentro de los gastos administrativos.

\subsection{Recursos humanos y cobertura del servicio}
\label{sec:e2-op-personas}

Cada perfil cuesta sus HH del mes por su valor en UF/HH, que corresponde a su categoría del Formulario E-26. La hoja \texttt{COSTOS} lleva la dotación en personas equivalentes y la pasa a HH con 160 HH por persona al mes. Entre los meses 21 y 56 se ocupan 107.778,05 HH, sin contar capacitación, que cuestan 74.273,19 UF. La Tabla~\ref{tab:e2-op-personas} separa mesa, NOC y SOC del resto del equipo.

\begin{tablalafrox}{Esfuerzo y costo de los perfiles en los meses 21 a 56. Fuente: modelo financiero, hoja COSTOS.}{tab:e2-op-personas}{Y R{2.0cm} R{2.0cm} R{2.9cm} R{2.2cm}}
  {\cab{Equipo} & \cab{HH} & \cab{UF} & \cab{CLP} & \cab{USD}}
  Mesa N1 & 38.246,00 & 25.624,82 & 1.024.992.800 & 1.138.881 \\
  NOC & 21.748,20 & 14.571,29 & 582.851.760 & 647.613 \\
  SOC & 26.280,00 & 17.607,60 & 704.304.000 & 782.560 \\
  Otros perfiles & 21.503,85 & 16.469,48 & 658.779.040 & 731.977 \\
  \midrule
  \textbf{Total sin capacitación} & \textbf{107.778,05} & \textbf{74.273,19} & \textbf{2.970.927.600} & \textbf{3.301.031} \\
\end{tablalafrox}

Las personas son el 68,11\,\% del costo de la Operación. Mesa N1, NOC y SOC suman 57.803,71 UF, y los demás perfiles 16.469,48 UF en coordinación, mantenimiento, integración, datos, calidad, seguridad y apoyo en terreno. Las horas de estos perfiles siguen las actividades del T-15 y no equivalen a puestos completos.

La mesa N1 atiende con puestos propios de 04:00 a 20:00 y suma un tercer agente en septiembre y diciembre desde el mes 25. De 20:00 a 22:00 la atiende el NOC, y en septiembre y diciembre la atención es de 24 horas. El NOC es compartido con la división 24×7 de LafroX y se carga con una dedicación de 100\,\% en el primer año de operación, 75\,\% en el segundo y 50\,\% en el tercero. El SOC parte en el mes 13, aunque aquí solo se cuentan sus horas desde el mes 21. Los tres perfiles cuestan 0,67 UF/HH y se cobran a 0,804 UF/HH (LafroX, 2026a). A esto se suman 1.933,53 UF de capacitación, mentoría y certificación y 314,21 UF de reservas de riesgo (LafroX, 2026c; 2026e).

\subsection{Infraestructura, alojamiento y servicios recurrentes}
\label{sec:e2-op-infraestructura}

La nube se costea con la estimación mensual de la arquitectura, las licencias por equipo y los enlaces por plan contratado. La Tabla~\ref{tab:e2-op-servicios} da el costo mensual de cada servicio y su total en los 36 meses.

\begin{tablalafrox}{Servicios recurrentes: costo mensual y total en los 36 meses de Operación. Fuente: modelo financiero, hojas ADQUISICIONES OPERACION, INNOVACIONES y COSTOS.}{tab:e2-op-servicios}{Y R{1.6cm} R{2.0cm} R{2.7cm} R{2.1cm}}
  {\cab{Servicio} & \cab{UF/mes} & \cab{UF total} & \cab{CLP total} & \cab{USD total}}
  Nube AWS & 368,92 & 13.280,97 & 531.238.855 & 590.265 \\
  Fibra empresarial & 13,60 & 489,60 & 19.584.000 & 21.760 \\
  LTE, Starlink y SIM & 41,85 & 1.506,63 & 60.265.107 & 66.961 \\
  Satélite Iridium SBD & 129,55 & 4.663,77 & 186.550.983 & 207.279 \\
  Protección de endpoints & 97,60 & 3.513,64 & 140.545.692 & 156.162 \\
  Microsoft Intune & 29,05 & 1.045,84 & 41.833.584 & 46.482 \\
  Custodia de respaldos & 4,96 & 178,57 & 7.142.857 & 7.937 \\
  Mapas, mensajería y pagos & 25,16 & 905,81 & 36.232.560 & 40.258 \\
  Nube de las innovaciones & Variable & 38,92 & 1.556.805 & 1.730 \\
  \midrule
  \textbf{Total proveedores} & \textbf{---} & \textbf{25.623,76} & \textbf{1.024.950.443} & \textbf{1.138.834} \\
\end{tablalafrox}

AWS cuesta 13.280,97 UF en la Operación, por contenedores, Aurora y DMS, almacenamiento y respaldos, red y VPN, seguridad y observabilidad, datos e IoT, analítica, acceso, recuperación ante desastres y entrega continua. De las 368,92 UF mensuales, 55,05 UF son soporte de AWS y reserva de consumo. La nube de las innovaciones agrega 38,92 UF, detalladas en la sección~\ref{sec:e2-inn-costos}.

Los equipos de la Operación cuestan 671,66 UF (\$26.866.393 o USD~29.852) y se pagan en el mes 21: 382,33 UF para reponer red y oficina, 257,33 UF de seguridad y 32,00 UF de adecuación eléctrica y puntos de red. Los equipos de terreno del CLIENTE quedan fuera del flujo, en la hoja \texttt{HARDWARE DEL CLIENTE}, y su instalación y mantención se cobran como horas de los perfiles de LafroX.

\subsection{Licencias, conectividad, energía e instalaciones}
\label{sec:e2-op-licencias}

CrowdStrike cuesta 19,99 USD por equipo al mes sobre 217 servidores y estaciones, 97,60 UF mensuales. Intune Plan 1 Device, con Claro Cloud Chile como proveedor, cuesta 3,38 USD por equipo para los 382 equipos Android del parque inicial, 29,05 UF al mes. Las dos licencias quedan a nombre del CLIENTE y su renovación está dentro de la cuota (Distribuidora Puelche S.A., 2026a, art.~84.3).

Los dos enlaces de fibra cuestan 6,80 UF cada uno al mes. Los ocho planes LTE, los cinco Starlink y hasta 198 SIM de terreno al año 3 suman 41,85 UF mensuales. El Iridium SBD de los 28 camiones con frío cuesta 129,55 UF al mes y envía temperaturas y alarmas en paralelo con la red móvil (LafroX, 2026a). Estas cuotas son solo el servicio, ya que los terminales se compran aparte. Mapas, mensajería y pagos agregan 25,16 UF al mes, y la custodia externa de respaldos 4,96 UF, que es distinta del depósito de códigos fuente de los gastos administrativos.

En energía e instalaciones, LafroX compra en la implementación la UPS, el generador, la climatización de precisión y las obras de la sala, paga 32,00 UF de adecuación eléctrica en el mes 21 y mantiene la sala con las horas de los paquetes 8.2.3 y 8.2.6. La electricidad que consumen la sala y los gabinetes sale del empalme de cada recinto del CLIENTE y no es costo de LafroX. Para estimarla, el Subdocumento 4 informa la carga eléctrica y el PUE de la sala (LafroX, 2026a; Distribuidora Puelche S.A., 2026b, RT-06.11).

\subsection{Seguros, garantías y costos de estructura}
\label{sec:e2-op-administrativos}

Los gastos administrativos de la Operación suman 6.225,30 UF (Tabla~\ref{tab:e2-op-administrativos}). Cada uno se carga en el mes en que se paga y no como un porcentaje sobre los costos directos.

\begin{tablalafrox}{Gastos administrativos de la Operación, meses 21 a 56. Fuente: modelo financiero, hoja GASTOS ADMINISTRATIVOS.}{tab:e2-op-administrativos}{Y R{2.2cm} R{3.0cm} R{2.4cm}}
  {\cab{Concepto} & \cab{UF} & \cab{CLP} & \cab{USD}}
  Overhead & 3.888,00 & 155.520.000 & 172.800 \\
  Garantías & 925,96 & 37.038.548 & 41.154 \\
  Seguros & 1.306,67 & 52.266.667 & 58.074 \\
  Custodia de fuentes & 104,67 & 4.186.875 & 4.652 \\
  \midrule
  \textbf{Total administrativo} & \textbf{6.225,30} & \textbf{249.012.090} & \textbf{276.680} \\
\end{tablalafrox}

El overhead es de 108 UF al mes: 37,50 UF de gerencia general, 45,00 UF de administración y finanzas, 15,00 UF de contabilidad y 10,50 UF de asesoría legal. Los gastos notariales se pagaron en el mes 1.

En garantías se pagan 825,00 UF de comisiones por la boleta de fiel cumplimiento y 100,96 UF por la de correcto funcionamiento. Se registra la comisión, de 1\,\% anual, y no el monto garantizado. La boleta de correcto funcionamiento equivale al 5\,\% del precio de implementación y dura dos años (Distribuidora Puelche S.A., 2026a, arts.~35 a 37).

Los seguros cuestan 1.306,67 UF, 426,67 UF de responsabilidad civil y 880,00 UF de ciberriesgo, con primas de 0,8\,\% y 1,2\,\% anual sobre límites de 20.000 UF. El seguro de equipos se pagó en la implementación (Distribuidora Puelche S.A., 2026a, art.~38). Las cuotas anuales del depósito de códigos fuente suman 104,67 UF.

\subsection{Seguimiento y optimizaciones durante la Operación}
\label{sec:e2-op-optimizacion}

Cada mes, el informe de nube del paquete 8.4.2 compara el consumo real con lo presupuestado, y cada trimestre el paquete 8.4.1 proyecta la capacidad que se necesitará en infraestructura, licencias y enlaces. Los cambios se proponen al CLIENTE antes de hacerlos (LafroX, 2026c, cuenta 8.4).

En la nube se ajusta el tamaño de contenedores y bases de datos, se apagan fuera de horario los ambientes que no son de producción y se revisan los plazos de retención, sin tocar los niveles de servicio ni el respaldo. En licencias y enlaces se cruzan altas y bajas con el inventario. En terreno, el mantenimiento preventivo y el plan de ciclo de vida del paquete 8.2.7 evitan reposiciones de urgencia.

El presupuesto no descuenta ahorros por estas medidas. Un ahorro se refleja recién cuando el CLIENTE aprueba el cambio que lo produce. La reducción de deuda técnica y la mejora continua usan horas ya programadas, cuyo costo se ve en la sección~\ref{sec:e2-mantencion}.
